% Section 3.3, from lectures/L03/appendix/b_exercises.md.

\section{Övningar}
\label{c3:sec:exercises}\label{c3:app:b}

\begin{aside}
\textbf{Så kontrollerar du ditt arbete.} Varje uttryck du läser av ur ett nät kan kontrolleras med
nätet självt: välj en rad i sanningstabellen, tänk dig knapparna i det läget, och följ strömvägen i
ritningen. Finns det en sluten väg från +24~V till lampan ska tabellen säga \code{1}.
Övning~\ref{c3:ex:3} kontrolleras på labbstationen.

Lösningsförslagen finns i bilaga~\ref{app:answers}.

Varje övning är märkt med sin sort. Övning~\ref{c3:ex:3} är kapitlets \textbf{kontrolluppgift}.
\end{aside}

\exercise{Läs av uttrycket}{Räkna för hand}
\label{c3:ex:1}

\bookfigure{l03_ex_networks}{Tre kontaktnät att analysera.}{c3:fig:ex_networks}

\noindent För vart och ett av näten a), b) och c):
\begin{parts}
  \item Skriv uttrycket för lampan, med parenteser runt varje grupp, enligt \secref{c3:a1}.
  \item Ta fram sanningstabellen.
  \item Sammanfatta tabellen med en mening.
\end{parts}

\exercise{Det minsta nätet}{Konstruktion}
\label{c3:ex:2}
\begin{parts}
  \item Förenkla vart och ett av de tre uttrycken från övning~\ref{c3:ex:1}, och ange vilken
    räknelag du använder.
  \item Beskriv det minsta nätet för varje uttryck: vilka kontakter, och hur de sitter.
  \item Hur många kontakter sparar du i varje nät? Hur många kontaktelement behöver varje knapp
    före och efter?
\end{parts}

\exercise{Kontroll: förutsäg och bygg}{Kontroll}
\label{c3:ex:3}
\emph{Förutsäg för hand, bygg på stationen, jämför.}

\bookfigure{l03_ex_kontroll}{Kontaktnätet i kontrolluppgiften.}{c3:fig:ex_kontroll}

\noindent Gör delarna i ordning, och skriv ned varje svar innan du går vidare.

\task{a) För hand}
Skriv uttrycket för H1, och förutsäg lampans läge för alla fyra kombinationerna av S1 och S2.

\task{b)}
Multiplicera ut uttrycket och förenkla. Vilken grind från kapitel~\ref{c2:ch} är det?

\task{c) På stationen}
Koppla nätet och pröva alla fyra kombinationerna. Stämmer lampan med din förutsägelse?

\task{d)}
Jämför med nätet för samma funktion i Labb~1, \hyperref[lab1:u6]{uppgift~6}. Hur många kontakter
har de två näten? Varför kan två så olika nät göra samma sak?

\exercise{Syntes}{Konstruktion}
\label{c3:ex:4}
\begin{parts}
  \item Ta fram det minsta kontaktnätet med två knappar för sanningstabellen nedan:

    \begin{narrowtable}{@{}ccc@{}}
      \thead{S1} & \thead{S2} & \thead{H1} \\
      \midrule
      0 & 0 & 1 \\
      0 & 1 & 1 \\
      1 & 0 & 0 \\
      1 & 1 & 1 \\
    \end{narrowtable}
  \item Ett transportband drivs av en motor som styrs av reläet K1. Motorn får gå om skyddsgrinden
    är stängd (givaren S3 påverkad), och samtidigt antingen startknappen på plats (S1) eller
    fjärrstarten (S2) är intryckt. Skriv uttrycket för K1 och rita strömvägen.
\end{parts}

\exercise{Felsökning}{Förståelse}
\label{c3:ex:5}

\bookfigure{l03_ex_fault}{Nätet som ska felsökas.}{c3:fig:ex_fault}

\noindent Nätet ska realisera \code{H1 = (S1 + S2) · S3'}.
\begin{parts}
  \item Du trycker på S1 men inte på S3, och lampan förblir släckt. Du mäter 24~V mellan 0~V och
    punkten ovanför parallellkopplingen, 24~V mellan 0~V och punkten mellan parallellkopplingen och
    S3, och 0~V mellan 0~V och punkten mellan S3 och lampan. Var sitter felet, och vad är den
    troligaste orsaken?
  \item Efter en ny koppling lyser lampan när S1 eller S2 är nedtryckt, men den slocknar inte när
    du dessutom trycker på S3. Vad är fel?
  \item Efter ytterligare en ny koppling lyser lampan bara när S3 \textbf{och} minst en av S1 och S2
    är nedtryckta. Vad är fel nu?
\end{parts}

\exercise{Stoppknappen}{Förståelse}
\label{c3:ex:6}
\begin{parts}
  \item I en maskins stoppkrets används stoppknappens brytande kontakt. Beskriv vad som händer med
    maskinen om ledningen till stoppknappen går av.
  \item Beskriv vad som skulle hända om stoppknappens slutande kontakt i stället användes, och
    kretsen var gjord så att maskinen stannar när kontakten sluts. Vad händer då om ledningen går
    av?
  \item Formulera en regel för hur en stoppfunktion ska kopplas, med ordet \emph{säker}.
\end{parts}

\exercise{Alla funktioner av två knappar}{Förståelse}
\label{c3:ex:7}
\begin{parts}
  \item En sanningstabell med två ingångar har fyra rader. Hur många olika sanningstabeller finns
    det, alltså hur många olika funktioner av två variabler?
  \item Två av funktionerna behöver inte någon kontakt alls. Vilka, och hur realiseras de?
  \item Vilka funktioner behöver både det slutande och det brytande kontaktelementet i båda
    knapparna?
\end{parts}

\exercise{En väljare av kontakter}{Räkna för hand, Fördjupning}
\label{c3:ex:8}
Ett nät har kontakterna S1 och S2 slutande parallellt, i serie med S1 brytande och S3 slutande
parallellt: \code{H1 = (S1 + S2) · (S1' + S3)}.
\begin{parts}
  \item Ta fram sanningstabellen.
  \item Dela tabellen i två halvor efter S1, och beskriv med ord vad nätet gör.
  \item Visa algebraiskt att \code{(S1 + S2)(S1' + S3) = S1' · S2 + S1 · S3}. Tänk på att
    \code{S1 · S1' = 0}, och att en term som \code{S2 · S3} kan visa sig vara onödig: kontrollera
    med tabellen.
  \item Jämför med grindnätet i \secref{c2:a8}. Vad har de gemensamt?
\end{parts}
